\documentclass[10pt,a4paper]{amsart}
\usepackage[margin=19mm]{geometry}
\usepackage{amsmath,amssymb,mathtools}
\usepackage[scr=rsfs,cal=euler]{mathalfa}
\usepackage{xfrac}
\usepackage[unicode,hidelinks,bookmarks=false]{hyperref}
\newcommand{\ie}{i.e.\ }
\newcommand{\cf}{cf.\ }
\newcommand{\resp}{respectively\ }
\newcommand{\Subsection}{\subsection}
\providecommand{\nobreakdash}{\nobreak\hbox{-}}
\setlength{\emergencystretch}{2em}
\multlinegap0pt
\newtheorem{thm}{Theorem}
\newtheorem{lem}[thm]{Lemma}
\newtheorem{defin}[thm]{Definition}
\newcommand{\IR}{{\mathbb R}}
\newcommand{\IN}{{\mathbb N}}
\newcommand{\IZ}{{\mathbb Z}}
\newcommand{\bx}{{\mathbf x}}
\newcommand{\bn}{{\mathbf n}}
\newcommand{\ri}{{\mathrm i}}
\newcommand{\re}{{\mathrm e}}
\newcommand{\rd}{{\,\mathrm d}}
\newcommand{\inc}{{\mathrm{inc}}}
\newcommand{\scat}{{\mathrm{scat}}}
\newcommand{\dn}{{\partial_\bn}}
\newcommand{\dxo}{{\partial_{x_1}}}
\newcommand{\dxt}{{\partial_{x_2}}}
\newcommand{\cu}{{\overline{u}}}
\newcommand{\cv}{{\overline{v}}}
\newcommand{\cw}{{\overline{w}}}
\newcommand{\GD}{{\Gamma_D}}
\newcommand{\GH}{{\Gamma_H}}
\newcommand{\GmH}{{\Gamma_{-H}}}
\newcommand{\GpmH}{{\Gamma_{\pm H}}}
\title[Explicit stability estimate for the periodic Helmholtz problem (Theorem 3.5)]{Trefftz Discontinuous Galerkin methods for scattering by periodic structures: the explicit stability estimate away from Rayleigh--Wood anomalies (Theorem 3.5)}
\author{Armando Maria Monforte, Andrea Moiola}
\date{Source: arXiv:2505.23216v2 (CC BY 4.0)}
\begin{document}
\begin{abstract}
	We propose a Trefftz discontinuous Galerkin (TDG) method for the approximation of plane wave scattering by periodic diffraction gratings, modelled by the two-dimensional Helmholtz equation.
	The periodic obstacle may include penetrable and impenetrable regions.
	The TDG method requires the approximation of the Dirichlet-to-Neumann (DtN) operator on the periodic cell faces, and relies on plane wave discrete spaces.
	For polygonal meshes, all linear-system entries can be computed analytically.
	Using a Rellich identity, we prove a new explicit stability estimate for the Helmholtz solution, which is robust in the small material jump limit.


\end{abstract}
\maketitle

\noindent\emph{Source and licence.} Trefftz Discontinuous Galerkin methods for scattering by periodic structures. Version arXiv:2505.23216v2 (CC BY 4.0), distributed under the Creative Commons Attribution 4.0 International licence, as recorded in the arXiv licence metadata for this exact version and shown on the abstract page https://arxiv.org/abs/2505.23216v2. Full rights notice: licenses/arxiv-2505.23216-CC-BY-4.0.txt.\par\smallskip

\noindent\textit{Editorial scope.} Counted unit: the theorem printed as Theorem 3.5 of the article (source0.tex lines 883--890, source label \texttt{thm-stima}), delivered together with the article's own complete original proof of it (source0.tex lines 891--1101, one \texttt{proof} environment of 211 lines that proves the whole explicit estimate with no gap and no deferral). No mathematics is written, repaired or re-derived: statement and proof are the article's own text line for line, in the article's own notation ($u$, $u^\inc$, $u^\scat$, $T^\pm$, $T^\pm\_M$, $\beta\_n^\pm$, $\alpha\_n$, $\Omega\_H^\pm$, $\Gamma\_{\pm H}$, $\delta$, $\kappa^\pm$, $H$, $L$, $\theta$). Required context retained uncounted: the article's abstract (lines 107--111), which announces exactly this stability estimate; its derivation of the outgoing Fourier representation of the scattered field together with the coefficients $\beta\_n^+$, the normal derivative (eq:dusdn) and the definition of the Dirichlet-to-Neumann operator $T^+$ (lines 442--468); its derivation of the coefficients $\beta\_n^-$ for real and for complex permittivity together with the definition of $T^-$ (lines 476--502); its definition of the truncated Dirichlet-to-Neumann operator $T^\pm\_M$ (lines 504--512, printed as Definition 2.4 of the article) and the article's Lemma 2.5 with the real and imaginary parts of the DtN pairings, including the display (eq:ImT)--(eq:ReT) that the counted proof uses (lines 514--527); the non-truncated boundary value problem (eq:nontrunc) with its incident wave (lines 621--632); the variational formulation (eq:weak) with the bilinear form (eq:bilinear) and the right-hand side (eq:right-hand-side) (lines 647--662); the interface notation $\Sigma$ and $\Gamma\_{j,j'}$ used by the Rellich identity (lines 673--684); the article's Lemma 3.1, the Rellich identity (eq:rellich), in statement form (lines 685--698); the article's Theorem 3.3, the non-trapping well-posedness theorem with its assumption (eq:NonTrap), in statement form (lines 792--805); the derivation of the coefficient identity (eq:Beta\_nExpansion) (lines 813--822); the display (eq:distance) defining the distance $\delta$ from the closest Rayleigh--Wood anomaly (lines 866--872); and the article's own paragraph describing the strategy of the counted proof (lines 875--882). Every definition, numbered display and label of those lines is retained unchanged. Omitted from this package and not redistributed: the article's preamble, title page, keywords and MSC block (lines 1--106); its whole introduction and the first three subsections of its second section with the quasi-periodic function spaces, the Fourier expansion and the scattering geometry (lines 112--441); the sentence introducing the continuity lemma of Pinto (2020) and that lemma (lines 469--475); the proofs of the retained Lemma 2.5 and the following Proposition 2.6 and Proposition 2.7 with the remainder of the DtN subsection (lines 528--618); the truncated boundary value problem (eq:trunc) and the paragraph introducing it (lines 633--644); the introduction of the variational formulation with its citation (lines 645--646); the opening of the stability section (lines 663--672); the proof of the Rellich identity and the whole subsection on solution existence and uniqueness including the article's Theorem 3.2 (lines 699--791); the proof of the retained Theorem 3.3 and the beginning of its argument (lines 806--812); the rest of the well-posedness proof and the opening of the subsection on explicit stability estimates before the display (eq:distance) (lines 823--865); the sentence pointing to the literature on DtN operators for acoustic waveguides (lines 873--874); the article's Remark 3.6 comparing its estimate with the literature, its whole fourth section on the DtN-TDG method with the formulation, the error analysis and the plane-wave assembly, its fifth section with the numerical experiments (lines 1102--2002); its bibliography (lines 2003--2004); and its closing end-document line (line 2005). Nothing inside a retained excerpt is omitted or altered: every retained body is delivered exactly as the article's lines are written, so no omission receipt of any kind accompanies this package. Citations: the retained excerpts carry exactly one citation command, in the optional argument of the retained Theorem 3.3, and it cites exactly one key of the article's own bibliography, \texttt{BonnetBD1994}. That entry is transcribed from the article's own bibliography, which is shipped with the source as the biber auxiliary file PeriodicTDG.bbl in the article's own entry list, and it is delivered with the article's own printed label, so the delivered citation mark reads [6] as the article's does (the printed reference list of the version PDF shows the same entry as number [6]). No other entry of the article's bibliography is transcribed or redistributed. Reference adaptation: none was needed and none was made. Every reference inside the delivered unit resolves inside it: the labels eq:nontrunc, eq:weak and eq:NonTrap of the retained displays, the labels eq:ImT, eq:ReT, eq:rellich, eq:beta\_piu, eq:beta\_meno, eq:distance and lemma:l2prod of the retained definitions, lemmas and displays, the label wellposed of the retained Theorem 3.3, and the labels eq:RellichIII, unmeno, unpiu, unresto, stimaorizzontale, stimal2, stimaverticale and eq:ReIGamma of the numbered displays inside the counted proof, together with the label thm-stima of the counted statement and its display (eq:stimaH1). No unresolved reference or citation is produced. Printed slips kept literal: no printed word, symbol, hypothesis or number of the counted statement or of its proof was altered; in particular the retained proof keeps the article's own repeated word in the sentence beginning ``In the last inequality we have have bounded''. Layout and class: the article's own class is the standard LaTeX \texttt{article} class at 10pt on a4paper with the packages babel, lmodern, latexsym, amsmath, amsthm, amssymb, mathtools, eucal, stmaryrd, graphicx, tikz, pgfplots, enumerate, biblatex, csquotes, hyperref and microtype; no publisher class or style file is used or shipped, so none travels with this package and none is redistributed. The wrapper is the lane's pinned \texttt{amsart} editorial wrapper at 10pt on A4, which adds the article's own macro definitions that the retained text needs ($\IR$, $\IN$, $\IZ$, $\bx$, $\bn$, $\ri$, $\re$, $\rd$, $\inc$, $\scat$, $\dn$, $\dxo$, $\dxt$, $\cu$, $\cv$, $\cw$, $\GD$, $\GH$, $\GmH$, $\GpmH$), declares the theorem-like environments the retained text needs (\texttt{thm}, \texttt{lem}, \texttt{defin}) with the article's own shared counter and without the article's section prefix, and keeps the standard environments the retained displays use (\texttt{equation}, \texttt{align}, \texttt{cases}). The delivered numbering is the unit's own: the retained Definition 2.4 prints as Definition 1, the retained Lemma 2.5 as Lemma 2, the retained Lemma 3.1 (Rellich identity) as Lemma 3, the retained Theorem 3.3 as Theorem 4, and the counted Theorem 3.5 as Theorem 5. The article's 10pt a4paper measure coincides with the wrapper's 10pt A4 measure, so no page-scale change of the mathematics is introduced. Assets: no retained span contains a figure environment, a graphics-inclusion command, a PDF-page inclusion, a table or a TikZ picture; the article's figures, graphics and TikZ drawings all live in the omitted parts of the paper, none is copied into this package, the package declares no asset dependency, and no image file is redistributed with it. Licence: version arXiv:2505.23216v2 of this article is distributed under the Creative Commons Attribution 4.0 International licence, as recorded in the arXiv licence metadata for this exact version and shown on the abstract page https://arxiv.org/abs/2505.23216v2. A full rights notice is delivered at licenses/arxiv-2505.23216-CC-BY-4.0.txt. Characters: the retained excerpts contain exactly two characters outside ASCII, both the accented letter of the word Poincar\'e as the article's own source writes it, one in the retained strategy paragraph and one inside the counted proof; the delivered engine is XeLaTeX, whose default Latin Modern text font carries that character, and the compile receipt reports no missing character. No glyph is substituted and no word is rewritten.
With these choices, the scattered field has the form
\begin{equation} \label{representation}
u^\scat(\bx)=\sum_{n\in\mathbb{Z}} u^\scat_n(H)\re^{\ri\beta_n^+ (x_2-H)} \re^{\ri\alpha_n x_1},    \qquad \bx\in\Omega_H^+ \cup \GH,
\end{equation}
where
\begin{equation} \label{beta_piu}
\beta_n^+ := \begin{cases} \sqrt{k^2\varepsilon^+-\alpha_n^2} & \alpha_n^2 \leq k^2\varepsilon^+, \\ 
	\ri\sqrt{\alpha_n^2-k^2\varepsilon^+} & \alpha_n^2 > k^2\varepsilon^+. \end{cases} 
	\end{equation}
	In particular, $\beta_0^+=-\kappa^+\sin\theta$ is real and positive.
	The normal derivative on $\GH$ of $u^\scat$ in \eqref{representation} is
	\begin{equation}\label{eq:dusdn}
\dxt u^\scat (x_1)
= \ri\sum_{n\in\mathbb{Z}} u^\scat_n(H) \beta_n^+ \re^{\ri\alpha_nx_1}.
\end{equation}
We want to enforce the relation between the value \eqref{representation} of $u^\scat$ and that of its normal derivative \eqref{eq:dusdn} in the formulation of the scattering problem,
so we define the Dirichlet-to-Neumann (DtN) operator $T^+$ as
\begin{align} \label{DtN}
& T^+:H^{1/2}_{\alpha_0}(\GH) \to H^{-1/2}_{\alpha_0}(\GH), \\ \nonumber
& (T^+\phi)(x_1) := \ri \sum_{n\in\mathbb{Z}} \phi_n \beta_n^+\re^{\ri\alpha_nx_1}, 
\hspace{0.5cm} \text{for} \hspace{0.2cm} 
\phi(x_1)=\sum_{n\in\IZ}\phi_n \re^{\ri\alpha_nx_1} \in H^{1/2}(\GH).
\end{align}

We say that a Helmholtz solution $u^\scat$ in $\Omega_H^+$ propagates upwards---equivalently, that it satisfies the radiation condition in $\Omega_H^+$---if 
$\gamma_H(\dxt u^\scat)=T^+(\gamma_H u^\scat)$, where $\gamma_H$ is the trace on $\GH$. 
In this case, $u^\scat$ admits the expansion \eqref{representation}.

We derive the expression of the DtN operator $T^-$ on $\GmH$.
In this case, we require the total field $u$---without subtracting the incoming wave $u^\inc$---to propagate downwards. 
Using a Fourier expansion and imposing that the total field solves the Helmholtz equation in $\Omega_H^-$ with $\varepsilon=\varepsilon^-$, we look for $u$ in the form
\begin{equation*}% \label{representation}
u(\bx)=\sum_{n\in\mathbb{Z}} u_n(-H)\re^{-\ri\beta_n^- (x_2+H)} \re^{\ri\alpha_n x_1},    \qquad \bx\in\Omega_H^-\cup\GmH.
\end{equation*}
Selecting the solutions similarly to (i)--(iii) above, if $\varepsilon^->0$ we obtain the coefficients
$\beta_n^-$ as
\begin{equation} \label{beta_meno}
\beta_n^- := \begin{cases} \sqrt{k^2\varepsilon^- - \alpha_n^2} & \alpha_n^2 \leq k^2\varepsilon^-, \\ 
	\ri\sqrt{\alpha_n^2-k^2\varepsilon^-} & \alpha_n^2 > k^2\varepsilon^-. \end{cases} 
	\end{equation}
	Instead, if $\varepsilon^-\notin\IR$, then the lower region $\Omega_H^-$ contains an absorbing medium, so all outgoing solutions decay exponentially for $x_2\to-\infty$.
	In this case we set 
	\begin{equation}\label{eq:betaComplEps}
\beta_n^-:=\sqrt{k^2\varepsilon^--\alpha_n^2} 
\end{equation}
choosing the complex root with positive imaginary part, i.e.\ $\Im\beta_n^->0$.
Since $k>0,\alpha_n\in\IR$ and $\Im\varepsilon^->0$, this is the standard branch cut of the square root, and we also have $\Re\beta_n^->0$.
Then, the DtN operator $T^-$ is
\begin{align} \label{DtN_meno}
& T^-:H^{1/2}_{\alpha_0}(\GmH) \to H^{-1/2}_{\alpha_0}(\GmH), \\ \nonumber
& (T^- \phi)(x_1) := \ri \sum_{n\in\mathbb{Z}} \phi_n \beta_n^- \re^{\ri\alpha_n x_1}, \hspace{0.5cm} \text{for} \hspace{0.2cm} 
\phi(x_1)=\sum_{n\in\IZ}\phi_n  \re^{\ri\alpha_nx_1} \in H^{1/2}_{\alpha_0}(\GmH),
\end{align}
and enjoys the same continuity property of $T^+$.
Note that the parameters $\alpha_n=k\sqrt{\varepsilon^+}\cos\theta+\frac{2\pi n}L$ enter the definition of the DtN operator $T^-$ on the lower boundary $\GmH$, but they depend on the value $\varepsilon^+$ of the material parameter $\varepsilon$ in the upper region $\Omega_H^+$, while $\varepsilon^-$ enters $T^-$ via the $\beta_n^-$ coefficients \eqref{beta_meno}--\eqref{eq:betaComplEps}.

In practical computations, we need to truncate the infinite series in the definitions of the DtN operators (\ref{DtN}) and (\ref{DtN_meno}), so we introduce the following operators

\begin{defin}[Truncated DtN operator]\label{def:TruncatedDtN}
For $M\in\IN$, we set
\begin{equation} \label{DtN_trunc}
	(T_M^\pm \phi)(x_1) := \ri \sum_{n=-M}^M \phi_n \beta_n^\pm \re^{\ri\alpha_nx_1},
	\qquad\forall \phi(x_1)=\sum_{n\in\IZ}\phi_n\re^{\ri\alpha_nx_1} \in H^{1/2}_{\alpha_0}(\GpmH).
\end{equation}
\end{defin}

\begin{lem} \label{lemma:l2prod}
Assuming $\varepsilon$ is real and positive, for every $w \in H^1_{\alpha_0} (\Omega)$,
\begin{align}\label{eq:ImT}
	\Im \int_\GpmH T^\pm w \, \cw  \rd s =
	L\sum_{\alpha^2_n < k^2\varepsilon^\pm} |w_n(\pm H)|^2 \sqrt{k^2\varepsilon^\pm-\alpha_n^2} \geq 0,\\ \label{eq:ReT}
	\Re \int_\GpmH T^\pm w \, \cw  \rd s =
	-L\sum_{\alpha^2_n > k^2\varepsilon^\pm} |w_n(\pm H)|^2 \sqrt{\alpha_n^2-k^2\varepsilon^\pm} \leq 0.
\end{align}
Moreover, for every $M \in \IN$,
\begin{align}
	\Im \int_\GpmH T^\pm w \, \cw  \rd s \geq \Im \int_\GpmH T_M^\pm w \, \cw  \rd s \geq 0,\\ 
	\Re \int_\GpmH T^\pm w \, \cw  \rd s \leq \Re \int_\GpmH T_M^\pm w \, \cw  \rd s \leq 0.
\end{align}
\end{lem}

Using the DtN operators $T^\pm$, we impose the appropriate boundary conditions on $\GpmH$, obtaining the following boundary value problem: given the incident wave 
$u^\inc(\bx) = \re^{\ri \kappa^+(x_1\cos \theta + x_2\sin \theta)} 
= \re^{\ri\alpha_0 x_1 - \ri \beta_0^+ x_2}$,
with $\theta \in \left[-\pi, 0 \right]$, find $u \in H^1_{\alpha_0}(\Omega)$, with $\alpha_0 = \kappa^+ \cos \theta$, such that
\begin{equation} \label{nontrunc}
\begin{cases}
	\Delta u + k^2\varepsilon u = 0 & \hspace{0.5cm} \text{in} \; \Omega, \\
	u = 0 & \hspace{0.5cm} \text{on} \; \GD, \\
	\dn(u-u^\inc) - T^+(u-u^\inc) = 0 & \hspace{0.5cm} \text{on} \; \GH, \\
	\dn u - T^-u = 0 & \hspace{0.5cm} \text{on} \; \GmH.
\end{cases}
\end{equation}

Let $u \in H^1_{\alpha_0, 0}(\Omega)$ be a distributional solution of the Helmholtz problem (\ref{nontrunc}). 
Multiplying both sides of the Helmholtz equation by a test function $v\in H^1_{\alpha_0, 0}(\Omega)$, and integrating by parts, the integrals over the left and right boundaries cancel by quasi-periodicity. Then, using the Dirichlet-to-Neumann operators to replace the normal derivative of the scattered field
$u^\scat=u-u^\inc$ on $\GpmH$, we obtain a weak formulation: find $u \in H^1_{\alpha_0, 0}(\Omega)$ such that
\begin{equation} \label{weak}
a_\varepsilon(u,v) = F(v) \hspace{0.5cm} \forall  v \in H^1_{\alpha_0, 0}(\Omega),
\end{equation}
where
\begin{align} \label{bilinear}
a_\varepsilon(u,v) & := 
\int_\Omega \left( \nabla u \cdot \nabla \cv - k^2\varepsilon u\cv \right) \rd\bx 
- \int_\GH  T^+u \; \cv \rd s 
- \int_\GmH T^-u \; \cv  \rd s, \\ \label{right-hand-side}
F(v) &  := -\int_\GH 2\ri\beta_0^+ u^\inc  \cv \rd s.
\end{align}
The value at the right-hand side \eqref{right-hand-side} comes from 
$-\dn u^\inc = T^+ u^\inc=\ri\beta_0^+ u^\inc$ on~$\GH$.


Suppose that $\Omega$ is divided into $P$ %different 
Lipschitz, connected subdomains $\Omega_j$, for $j=1,\ldots,P$, in which the relative permittivity $\varepsilon$ assumes a constant value.
Denote~by 
\begin{align*}
\Sigma:=\big\{&(j,j')\in \{1,\ldots,P\}^2: \; j<j',\\
&\Gamma_{j,j'}:=\partial\Omega_j\cap\partial\Omega_{j'} \text{ has positive 1-dimensional measure}\big\} 
\end{align*}
the index set of the interfaces between constant-$\varepsilon$ regions.
On $\Gamma_{j,j'}$ with $(j,j')\in\Sigma$, let $\bn=(n_1,n_2)$ denote the unit normal pointing from $\Omega_j$ into $\Omega_{j'}$.



\begin{lem}[Rellich identity] \label{RelLem}
Assume that $\varepsilon$ is real and positive. 
If $u\in H^1_{\alpha_0,0}(\Omega)$ is a solution to the variational problem (\ref{weak}) and the trace of $u$ on $\GpmH$ belongs to $H^1(\GpmH)$, then the following Rellich identity holds:
\begin{align} \nonumber
	& \int_\Omega 2 \left| \dxt u \right|^2 \rd\bx
	-  k^2\sum_{(j,j')\in\Sigma} ( \varepsilon_{j} -\varepsilon_{j'} ) \int_{\Gamma_{j,j'}} x_2 n_2 |u|^2\rd s
	- \int_{\GD} x_2n_2|\dn u|^2 \rd s\\
	&+ H \int_{\GH\cup\GmH} \biggl( \left|\dxo u \right|^2 - \left| \dn u \right|^2 -k^2\varepsilon |u|^2 \biggr)\rd s
	-\int_{\GH} T^+ u \cu\rd s -\int_\GmH T^- u \cu \rd s
	\nonumber \\
	& = -\int_{\GH} 2 \ri \beta_0^+ u^\inc \cu \rd s.
	\label{rellich}
\end{align}
\end{lem}

\begin{thm}[{\cite[Theorem~3.5]{BonnetBD1994}}] \label{wellposed} 
Assume that $\varepsilon$ is positive real and
\begin{equation}\label{eq:NonTrap}
	\begin{cases}
		n_2 \, x_2 \leq 0 \quad\text{on} \hspace{0.2cm} \GD, \\
		\forall x_1\in[0,L], \; \tau\mapsto\varepsilon(x_1,\tau)\text{ and }
		\tau\mapsto\varepsilon(x_1,-\tau)\\
		\hspace{20mm}\text{are monotonically non-decreasing for }\tau\in[0,\infty),
	\end{cases}
\end{equation}
where $\bn = (n_1, n_2)$ is the unit normal on $\GD$ pointing from $\Omega$ into $D$.
Assume that the scattering problem is non-trivial, i.e.\ $D\ne\emptyset$ or $\varepsilon$ is not constant (or both).
Then (\ref{nontrunc}) is well-posed for every value of $k$.
\end{thm}

$u(x_1,\pm H)=\sum_{n\in\IZ}u_n^\pm \re^{\ri\alpha_nx_1}$.
The imaginary part of the Rellich identity \eqref{rellich} with $u^\inc=0$, together with \eqref{eq:ImT}, gives that $u_n^\pm=0$ for all $n$ with $\alpha_n^2<k^2\varepsilon^\pm$.
The definitions \eqref{beta_piu} and \eqref{beta_meno} of $\beta_n^\pm$ give
\begin{equation}\label{Beta_nExpansion}
	\alpha_n^2-|\beta_n^\pm|^2-k^2\varepsilon^\pm =-(\beta_n^\pm)^2-|\beta_n^\pm|^2 =
	\begin{cases}
		-2|\beta_n^\pm|^2 &\alpha_n^2\le k^2\varepsilon^\pm,\\
		0&\alpha_n^2> k^2\varepsilon^\pm.
	\end{cases}
\end{equation}

This means that one of the two plane waves $\re^{\ri \kappa^\pm x_1}$, propagating in $\Omega^\pm$ in direction $x_1$, are quasi-periodic with parameter $\alpha_0$ over the interval $[0,L]$. 
We denote by $\delta$ a measure of the distance from the closest Rayleigh--Wood anomaly:
\begin{equation} \label{distance}
\delta_\pm := \min_{n\in\IZ} |\beta_n^\pm|, \qquad
\delta := \min\{ \delta_+, \delta_- \},
\end{equation}
and in the following analysis we assume that its value is nonzero.

The next theorem gives an explicit quantitative bound on the (wavenumber-weighted) $H^1(\Omega)$ norm of the solution under this assumption.
The key tool in its proof is the Rellich identity \eqref{rellich}.
The non-trapping condition, together with Lemma~\ref{lemma:l2prod}, determines the signs of the integrals in \eqref{rellich}.
The non-anomaly assumption $\delta>0$ and the integrals over $\GpmH$ in \eqref{rellich} allow to control all the terms in the Fourier expansion of $u$ on $\GpmH$ (see \eqref{unmeno}--\eqref{unpiu}).
From this, we bound $\|u\|_{L^2(\GH\cup\GmH)}$.
The norm in $\Omega$ is controlled using a Poincaré-type inequality in the $x_2$ direction, taking advantage of the $\|\partial_{x_2}u\|_{L^2(\Omega)}$ term in the Rellich identity.



\begin{thm} \label{thm-stima}
Under the non-trapping assumptions \eqref{eq:NonTrap}, and assuming that $\delta$ in \eqref{distance} is nonzero,  the unique solution $u$ of \eqref{nontrunc} satisfies the stability bound
\begin{align}\label{stimaH1} 
	\|\kappa u\|_{L^2(\Omega)}^2 + \|\nabla u\|_{L^2(\Omega)}^2
	\le 10L\kappa_+|\sin\theta|\bigg(\frac H\delta\|\kappa\|_{L^\infty(\Omega)}^2
	\big(2+7H\max\{\kappa^+,\kappa^-\}\big)^3+1\bigg).
\end{align}
\end{thm}

\begin{proof}
Let $u$ be the solution of \eqref{nontrunc}.
Under the non-trapping assumption \eqref{eq:NonTrap}, we know by Theorem~\ref{wellposed} that $u$ exists and it is uniquely defined.
We expand the outward-propagating components, 
i.e.\ $u^\scat$ in $\Omega_H^+$ and $u$ in $\Omega_H^-$, in Fourier series:
\begin{align*}
	&u(\bx)=
	\begin{cases}\displaystyle
		u^\inc(\bx)+u^\scat(\bx)
		=\re^{\ri(\alpha_0x_1-\beta_0^+x_2)}+
		\sum_{n \in \IZ}u_n^\scat\re^{\ri(\alpha_nx_1+\beta_n^+ (x_2-H))} & \bx\in\GH\cup \Omega_H^+,
		\\\displaystyle
		\sum_{n \in \IZ}u_n^-\re^{\ri(\alpha_nx_1-\beta_n^- (x_2+H))} & \bx\in\GmH\cup \Omega_H^-.
	\end{cases}
\end{align*}
In particular,
$u(\bx)=(u_0^\scat+\re^{-\ri\beta_0^+H})\re^{\ri\alpha_0x_1}
+\sum_{0\ne n \in \IZ}u_n^\scat\re^{\ri\alpha_nx_1}$ on $\GH$.
Their partial derivatives are:
\begin{align*}
	&\dxo u(\bx)=
	\begin{cases}\displaystyle
		\ri\alpha_0\re^{\ri(\alpha_0x_1-\beta_0^+x_2)}+
		\ri\sum_{n \in \IZ}\alpha_n u_n^\scat\re^{\ri(\alpha_nx_1+\beta_n^+ (x_2-H))} & \bx\in\GH\cup \Omega_H^+,
		\\\displaystyle
		\ri\sum_{n \in \IZ}\alpha_n u_n^-\re^{\ri(\alpha_nx_1-\beta_n^- (x_2+H))} & \bx\in\GmH\cup \Omega_H^-,
	\end{cases}
	\\
	&\dxt u(\bx)=
	\begin{cases}\displaystyle
		-\ri\beta_0^+\re^{\ri(\alpha_0x_1-\beta_0^+x_2)}+
		\ri\sum_{n \in \IZ}\beta_n^+ u_n^\scat\re^{\ri(\alpha_nx_1+\beta_n^+ (x_2-H))} & \bx\in\GH\cup \Omega_H^+,
		\\\displaystyle
		-\ri\sum_{n \in \IZ}\beta_n^- u_n^-\re^{\ri(\alpha_nx_1-\beta_n^- (x_2+H))} & \bx\in\GmH\cup \Omega_H^-.
	\end{cases}%\\
\end{align*}
Identity \eqref{Beta_nExpansion} allows to expand the $\GpmH$ term in the Rellich identity \eqref{rellich} in terms of the Fourier coefficients $u_n^\scat$ and $u_n^-$:
\begin{align*}
	&\int_{\GH\cup \GmH} \biggl(\left| \dxo u \right|^2 - \left| \dn u \right|^2 -k^2\varepsilon |u|^2 \biggr) \rd s \\
	&= L(\alpha_0^2-k^2\varepsilon^+)|\re^{-\ri \beta_0^+H}+u_0^\scat|^2
	-L|\beta_0^+|^2|\re^{-\ri \beta_0^+H}-u_0^\scat|^2\\
	&\qquad-2L \sum_{0\neq n \in \IZ, \alpha_n^2 \leq k^2\varepsilon^+} |\beta_n^+|^2 {|u_n^\scat|^2}
	- 2L \sum_{n \in \IZ, \alpha_n^2 \leq k^2\varepsilon^-} |\beta_n^-|^2|u^-_n|^2.
\end{align*}
Using the definitions \eqref{beta_piu}, \eqref{DtN}, \eqref{beta_meno}, \eqref{DtN_meno} of $T^\pm$ and $\beta_n^\pm$, we expand the DtN operators:
\begin{align*}
	&\int_{\GH} T^+ u\cu\rd s
	=-L\sum_{\alpha^2_n > k^2\varepsilon^+} |\beta_n^+||u_n^\scat|^2
	+\ri L |\beta_0^+||u_0^\scat+\re^{-\ri \beta_0^+H}|^2
	+\ri L\!\!\sum_{0\ne n\in\IZ,\;\alpha^2_n < k^2\varepsilon^-}\!\!\! |\beta_n^+||u_n^\scat|^2 ,
	\\
	&\int_\GmH T^- u\cu\rd s
	=-L\sum_{\alpha^2_n > k^2\varepsilon^-} |\beta_n^-||u_n^-|^2
	+\ri L\sum_{\alpha^2_n \leq k^2\varepsilon^-} |\beta_n^-||u_n^-|^2 .
\end{align*}

We separate the terms appearing in the Rellich identity \eqref{rellich} associated to $\GpmH$, to the other parts of the domain, and to the right-hand side:
\begin{align*} \nonumber
	I_\GpmH := & H\int_{\GH\cup\GmH} \biggl( \left|\dxo u \right|^2 - \left| \dn u \right|^2 -k^2\varepsilon |u|^2 \biggr)\rd s -\int_{\GH} T^+ u\cu\rd s-\int_\GmH T^- u\cu\rd s
	\\ \nonumber
	=& 
	-L\Big(H|\beta_0^+|^2+\ri|\beta_0^+|\Big) \big|\re^{-\ri\beta_0^+ H}+u^\scat_0\big|^2
	- L H|\beta_0^+|^2\big|\re^{-\ri\beta_0^+ H}-u^\scat_0\big|^2\\ \nonumber
	&-L\sum_{0\ne n\in\IZ,\; \alpha_n^2\le k^2\varepsilon^+}\Big(2H|\beta_n^+|^2+\ri|\beta_n^+|\Big) |u^{\scat}_n|^2
	-L\sum_{ \alpha_n^2\le k^2\varepsilon^-}\Big(2H|\beta_n^-|^2+\ri|\beta_n^-|\Big)|u^-_n|^2\\ 
	%\label{eq:rellich1}
	&+L\sum_{\alpha^2_n > k^2\varepsilon^+} |\beta_n^+||u_n^{\scat}|^2
	+L\sum_{\alpha^2_n > k^2\varepsilon^-}|\beta_n^-||u_n^-|^2,
	\\
	I_{\Omega, \Sigma, \GD} :=& \int_\Omega 2|\dxt u|^2\rd\bx
	-k^2\sum_{(j,j')\in\Sigma} ( \varepsilon_{j} -\varepsilon_{j'} ) \int_{\Gamma_{j,j'}} x_2 n_2 |u|^2\rd s
	- \int_{\GD} x_2n_2|\dn u|^2 \rd s,
	\\
	I_{\mathrm{RHS}} :=& -\int_{\GH} 2 \ri \beta_0^+ {u^\inc} \cu \rd s
	=-2\ri L\beta_0^+ {(1+\overline{u_0^\scat} \re^{-\ri\beta_0^+H})}.
\end{align*}
The Rellich identity writes as 
\begin{equation}\label{eq:RellichIII}
	I_\GpmH+I_{\Omega, \Sigma, \GD}=I_{\mathrm{RHS}}, \qquad
	\text{with} \quad I_{\Omega, \Sigma, \GD}\in \IR, \quad
	I_{\Omega, \Sigma, \GD}\ge 0, \quad\Im I_\GpmH\le0, 
\end{equation}
thanks to the non-trapping assumption \eqref{eq:NonTrap}, which implies that $(\varepsilon_j-\varepsilon_{j'})x_2n_2\le0$ on $\Gamma_{j,j'}$, $(j,j')\in\Sigma$.
For simplicity we use the coefficients of the total field $u$ on $\GH$:
$$
u_n^+:=\begin{cases} u_0^\scat+\re^{-\ri\beta_0^+H}  & n=0,\\ u_n^\scat & n\ne0.\end{cases}
$$
It follows that $I_{\mathrm{RHS}}=-2\ri L\beta_0^+ \overline{u_0^+}\re^{-\ri\beta_0^+H}$.
Thanks to \eqref{eq:RellichIII} and the weighted Young inequality,
\begin{align*}
	L\sum_{\alpha_n^2\le k^2\varepsilon^+}|\beta_n^+||u^+_n|^2
	+L\sum_{\alpha_n^2\le k^2\varepsilon^-}|\beta_n^-||u^-_n|^2
	& =|\Im I_\GpmH|
	=|\Im I_{\mathrm{RHS}}|
	\le 2L |\beta_0^+||u_0^+|\\
	&\le 2L|\beta_0^+|+\frac L2|\beta_0^+||u_0^+|^2,
\end{align*}
so, bringing the last term of the inequality to the left and dividing by $\frac12$, we can control the propagative-mode coefficients:
\begin{align} \label{unmeno}
	\sum_{\alpha_n^2\le k^2\varepsilon^+}|\beta_n^+||u^+_n|^2
	+\sum_{\alpha_n^2\le k^2\varepsilon^-}|\beta_n^-||u^-_n|^2
	&\le 4|\beta_0^+|,
	\\ \label{unpiu}
	\sum_{\alpha_n^2\le k^2\varepsilon^+}|u^+_n|^2
	+\sum_{\alpha_n^2\le k^2\varepsilon^-}|u^-_n|^2
	&\le \frac{4}{\delta}|\beta_0^+|.
\end{align}
In particular, \eqref{unmeno} gives $|u_0^+|\le2$ and $|I_{\mathrm{RHS}}|\le 4L|\beta_0^+|$.
Moreover,
\begin{align}\nonumber
	\Re I_\GpmH
	=&-LH|\beta_0^+|^2 |u^+_0|^2
	-2LH\sum_{0\ne n\in\IZ,\; \alpha_n^2\le k^2\varepsilon^+}|\beta_n^+|^2|u^+_n|^2 
	-2LH\sum_{n\in\IZ,\; \alpha_n^2\le k^2\varepsilon^-}|\beta_n^-|^2|u^-_n|^2  \\
	& - L H|\beta_0^+|^2\big|2\re^{-\ri\beta_0^+ H}-u^{+}_0\big|^2
	+L\sum_{\alpha^2_n > k^2\varepsilon^+} |\beta_n^+||u_n^+|^2
	+L\sum_{\alpha^2_n > k^2\varepsilon^-} |\beta_n^-||u_n^-|^2,
	\label{eq:ReIGamma}
\end{align}
so\, from $\Re I_\GpmH\le |I_{\mathrm{RHS}}| {\le  4L|\beta_0^+|}$,
\begin{align} \nonumber
	&\sum_{\alpha^2_n > k^2\varepsilon^+} |\beta_n^+||u_n^+|^2
	+\sum_{\alpha^2_n > k^2\varepsilon^-} |\beta_n^-||u_n^-|^2
	\\
	\nonumber
	\le	&  \, {4|\beta_0^+|} 
	+20H|\beta_0^+|^2 
	+2H\sum_{0\ne n\in\IZ,\; \alpha_n^2\le k^2\varepsilon^+}|\beta_n^+|^2|u^+_n|^2 
	+2H\sum_{ \alpha_n^2\le k^2\varepsilon^-}|\beta_n^-|^2|u^-_n|^2 
	\\ \label{unresto}
	\le&
	\big(4+28H\max\{\kappa^+,\kappa^-\}\big)|\beta_0^+|,
\end{align}
where we used 
inequality~\eqref{unmeno} and
$|\beta_n^\pm|\le\kappa^\pm$ 
for $n$ with $\alpha_n^2\le k^2\varepsilon^\pm$.
From \eqref{unpiu} and \eqref{unresto}, and recalling the definition~\eqref{distance} of $\delta$, we get an estimate for the norm of the solution on $\GH \cup \GmH$
\begin{align}%\nonumber
	\|u\|_{L^2(\GH\cup\GmH)}^2= L
	\sum_{n\in\IZ}(|u^+_n|^2+ |u^-_n|^2)
	&\le
	{\frac{4L}\delta\big(2+7 H\max\{\kappa^+,\kappa^-\}\big)|\beta_0^+|. }
	\label{stimaorizzontale}
\end{align}
We define 
\begin{equation*}
	\Omega_+:=\Omega\cap\{x_2>0\}, \quad
	\Omega_-:=\Omega\cap\{x_2<0\}, \quad
	m_+(x_1):=\inf\{x_2\ge 0: (x_1,x_2)\in\Omega\}, \; 0<x_1<L,
\end{equation*}
and use 
a Poincaré-type inequality 
in the $x_2$ variable 
to bound the $L^2$ norm on $\Omega_\pm$
\begin{align*}
	\|u\|_{L^2(\Omega_+)}^2 & =
	\int_0^L\int_{m_+(x_1)}^H \bigg|u(x_1,H)-\int_{x_2}^H \dxt u(x_1,x')\rd x'\bigg|^2\rd x_2\rd x_1\\&
	\le 2H\|u\|_{L^2(\GH)}^2 + H^2\|\dxt u\|_{L^2(\Omega_+)}^2.
\end{align*}
The same holds for $\Omega_-$ and $\GmH$, so 
\begin{equation} \label{stimal2}
	\|u\|_{L^2(\Omega)}^2\le 2H\|u\|_{L^2(\GH\cup\GmH)}^2 + H^2\|\dxt u\|_{L^2(\Omega)}^2.
\end{equation}
Recalling \eqref{eq:RellichIII}, $|I_{\mathrm{RHS}}|\le 4L|\beta_0^+|$, and collecting the negative terms in \eqref{eq:ReIGamma},
\begin{align} \nonumber
	2\|\partial_{x_2}u\|^2_{L^2(\Omega)}
	\le&I_{\Omega, \Sigma, \GD}
	=\Re I_{\Omega, \Sigma, \GD}
	=\Re I_{\mathrm{RHS}}-\Re I_\GpmH
	\\ \nonumber
	\le &  
	4L|\beta_0^+| + 20 LH|\beta_0^+|^2 
	+2LH\!\!\sum_{0\ne n\in\IZ,\; \alpha_n^2\le k^2\varepsilon^+}|\beta_n^+|^2|u^+_n|^2 
	+2LH\sum_{\alpha_n^2\le k^2\varepsilon^-}|\beta_n^-|^2|u^-_n|^2 
	\\ \label{stimaverticale}
	\le& 4L|\beta_0^+| + 20 LH|\beta_0^+|^2 +
	2 k^2 H \Big(\varepsilon^+\|u\|_{L^2(\GH)}^2+\varepsilon^-\|u\|_{L^2(\GmH)}^2\Big)
	,\end{align}
where we used that $|\beta_n^\pm|^2\le k^2\varepsilon^\pm$ for these $n$.
From \eqref{stimal2}, \eqref{stimaorizzontale}, \eqref{stimaverticale}, and again $\beta_0^+\le\kappa^+$, we control $u$ over the whole domain:
\begin{align*} \nonumber
	&\|u\|_{L^2(\Omega)}^2\\ \nonumber
	&\le   2H\|u\|_{L^2(\GH\cup\GmH)}^2 
	+2LH^2|\beta_0^+| + 10 LH^3|\beta_0^+|^2 +
	k^2 H^3 \Big(\varepsilon^+\|u\|_{L^2(\GH)}^2+\varepsilon^-\|u\|_{L^2(\GmH)}^2\Big)
	\\ \nonumber
	&\le \frac{4LH}\delta\big(2+H^2\max\{\kappa^+,\kappa^-\}^2\big)\big(2+7H\max\{\kappa^+,\kappa^-\}\big)|\beta_0^+|
	+2LH^2(1+5H\kappa^+)|\beta_0^+|
	\\
	&\le \frac{5LH}\delta\big(2+7H\max\{\kappa^+,\kappa^-\}\big)^3|\beta_0^+|.
\end{align*}
In the last inequality we have have bounded the last term with a quarter of the first one, using $2H\le\frac1\delta+H^2\delta\le\frac1\delta(1+H^2(\kappa^+)^2)$, which follows from $\delta\le\kappa^+$.
To conclude, we bound the gradient of $u$ using the weak formulation \eqref{weak}, the sign \eqref{eq:ReT} of $\Re T^\pm$, $|I_{\mathrm{RHS}}|\le 4L\beta_0^+$, and $\beta_0^+=-\kappa^+\sin\theta>0$:
\begin{align*} 
	\|u\|_{\kappa, H^1(\Omega)}^2 = 
	&\|\kappa u\|_{L^2(\Omega)}^2 + \|\nabla u\|_{L^2(\Omega)}^2
	\\ 
	=&\|\kappa u\|_{L^2(\Omega)}^2 +
	\int_\Omega \kappa^2 |u|^2\rd\bx+
	\Re\Big\{\int_\GH T^+u\cu\rd s+\int_\GmH T^-u\cu\rd s\Big\}+\Re I_{\mathrm{RHS}}
	\\ 
	\le& 2 \|\kappa u\|_{L^2(\Omega)}^2 +4L|\beta_0^+|
	\\ 
	\le& \frac{10LH}\delta\|\kappa\|_{L^\infty(\Omega)}^2 \big(2+7H\max\{\kappa^+,\kappa^-\}\big)^3|\beta_0^+|
	+4L|\beta_0^+|
	\\ 
	\le& 10L\kappa_+|\sin\theta|\bigg(\frac H\delta\|\kappa\|_{L^\infty(\Omega)}^2
	\big(2+7H\max\{\kappa^+,\kappa^-\}\big)^3+1\bigg).
\end{align*}
\end{proof}

\begin{thebibliography}{99}
\bibitem[6]{BonnetBD1994} A.-S. Bonnet-Bendhia, F. Starling, \emph{Guided waves by electromagnetic gratings and nonuniqueness examples for the diffraction problem}, Math. Methods Appl. Sci. \textbf{17} (1994), no.~5, 305--338.
\end{thebibliography}
\end{document}
